\item \textbf{Problema de repaso.} Después de una colisión en el espacio exterior, un disco de cobre a 850 °C gira en torno a su eje con una rapidez angular de 25.0 rad/s. A medida que el disco radia luz infrarroja, su temperatura cae a 20.0 °C. Ningún momento de torsión externo actúa en el disco. (a) ¿La rapidez angular cambia a medida que el disco se enfría? Explique cómo cambia o por qué no lo hace. (b) ¿Cuál es su rapidez angular a la temperatura más baja?
